Let \((\xi,\gamma)\in\mathcal F(\kappa)\). By definition, \(\xi\gamma>d\) and there is a law \(P\in\mathcal M(\xi,\gamma;\kappa)\). The compatibility lemma gives \(\xi=1\) and \(\gamma>d\). Thus
\[
\mathcal F(\kappa)\subseteq\{(1,\gamma):\gamma>d\}.
\]
Conversely, \(\mathrm{thm:calibrated-witness}\) supplies a member of \(\mathcal M(1,\gamma;\kappa)\) for every \(\gamma>d\), so the reverse inclusion holds. This proves the frontier equality.

For every feasible pair, \(\xi=1\); therefore \(\mathrm{def:transverse-model}\) gives \(\mathcal M(\xi,\gamma;\kappa)=\mathcal M_{\mathrm{tr}}(\gamma;\kappa)\). Applying \(\mathrm{thm:transverse-minimax-rate}\) and observing that \(1+\xi=2\) gives
\[
c_1(r_n^*)^{1+\xi}\le\mathfrak R_n(\xi,\gamma;\kappa)\le C_1(r_n^*)^{1+\xi}.
\]
The same theorem identifies \(\widehat G_n^{\mathrm{LP}}\) as the attaining estimator in each branch. Hence the frozen calibrated problem has no feasible nontransverse super-margin labels: the apparent general frontier collapses exactly to the transverse line.